\section{Resultados}\label{sec:resultados}

% ─── Beat 1: Visão consolidada. ─────────────────────────────────────
\subsection{Visão Geral}
% Tabela 3-vias × 6 cenários × métricas principais.
% Headline: AIMD CUBIC ganha em pico de RTT em 5/6 cenários
% (maior ganho: -43\% em A) e em tempo total de stall em 5/6
% (maior ganho: 11 s → 4,3 s em E), mas eleva p50/p95 em todos
% e perde ratings subjetivos em 3/6.

% ─── Beat 2: Por cenário. ───────────────────────────────────────────
\subsection{Análise por Cenário}

\subsubsection{Cenário A — Queda Abrupta}
% Caso onde CUBIC mais brilha: -43\% no pico.

\subsubsection{Cenário B — Degradação em Degraus}

\subsubsection{Cenário C — Dois Patamares}

\subsubsection{Cenário D — Janelas Curtas}

\subsubsection{Cenário E — Constrição Pesada Longa}
% Headline secundário: stalls 11 s → 4,3 s.

\subsubsection{Cenário F — Regressão Discutida}
% Único cenário em que CUBIC regride. Discutir por quê.

% ─── Beat 3: Trade-off central. ─────────────────────────────────────
\subsection{O Trade-off: Cauda vs. Regime Estacionário}
% CUBIC otimiza cauda da distribuição (pico, p99, stalls);
% AIMD de base otimiza regime estacionário (p50, p95).
% Esta é uma escolha de prioridade, não dominância.
% Conexão com \cite{alos2019} §6 (design philosophy).

% ─── Beat 3.5: Análise de sensibilidade ao design (5b/5c). ─────────
\subsection{Análise de Sensibilidade ao Design da Curva}\label{sec:sensibilidade}
% Após a comparação principal, avaliamos duas variantes de design da
% curva cúbica para entender o porquê da forma final adotada:
%
% - Variante 5b: âncora temporal em \texttt{rtt\_clear\_since} em vez
%   de \texttt{last\_decrease\_at}.
% - Variante 5c: acumulação monotônica de t entre crescimentos
%   (CUBIC literal, §3.1 do paper original).
%
% Ambas regridem o pico de RTT em cenários constrangidos (A K=5 s,
% E K=15 s). Diagnóstico estrutural: o CUBIC foi originalmente
% formulado para amostragem por-ACK, regime em que pequenos deltas
% acumulam ao longo de muitas iterações; em cadência de 1 Hz, a curva
% aplica o delta cheio em cada tick, resultando em overshoot
% sistemático.
%
% Achado de pesquisa a destacar: o CUBIC literal pressupõe densidade
% de amostragem por-ACK; em ambientes de cadência reduzida (1 Hz no
% nosso caso), a variante simplificada (âncora em last_decrease_at,
% reset de clear_since a cada crescimento) é não apenas suficiente,
% mas superior.
%
% Fontes: [[results-idea-5b]], [[results-idea-5c]].

% ─── Beat 4: Ratings subjetivos. ────────────────────────────────────
\subsection{Avaliação Subjetiva}
% Tabela de ratings 1-5 por cenário e controlador.
% n=1 avaliador. Limitação reconhecida.

% ─── Beat 5: Caveats. ───────────────────────────────────────────────
\subsection{Limitações dos Resultados}
% n=1 execução por célula CUBIC; reaproveitamento das colunas
% base; conteúdo do jogo afeta dificuldade efetiva (encoder produz
% menos quando cena é menos demandante); um único avaliador subjetivo.
